
\sourcenote{Genre: admissible rule (tier 2) --- a proof-tree-rewriting
recipe. The rule is sequent-side machinery: premises and conclusion are
all $\vdash$ rows, and what the extension theory contributes is the
\emph{scope} vocabulary --- the child atom $\{g_q \ge 0\}$ names a
region concluded by a nested certificate, a shape no base cell can
express. The rule closes a named trust boundary: before its
mechanization, the move was performed producer-side and asserted.}

\begin{brecipe}[the admissible template, general form]\label{rec:hm-template}
An admissible rule $R$ with premises $P_1, \dots, P_n$ and conclusion
$C$ is warranted by a \emph{rewrite}: any proof-tree node applying $R$
is replaced by a derivation fragment over core (and
previously-justified) rules with the same conclusion, consuming the
premise subtrees at its leaves --- or, where no catalogue derivation of
$C$ can exist, by a discharge against a named verified theorem (the
degenerate rewrite, the tier\textquotesingle s recorded boundary case).
Leaf-first application eliminates every $R$-node from a finite tree in
finitely many steps, so soundness is conserved with no new kernel
semantic content.
\end{brecipe}

\begin{brecipe}[the \rulename{HypMultiplier} instance]\label{rec:hm-instance}
\[
\inferrule*[Right=\rulename{HypMultiplier}]{
  \mathcal{T}_\lambda \rhd\ \vdash \lambda \ge 0 \ \text{on } \mathrm{cell} \\
  \mathcal{T}_c \rhd\ \vdash N + \lambda\, g_q \le 0 \ \text{on } \mathrm{cell}
}{
  \vdash N \le 0 \ \text{on } \mathrm{cell} \cap \{g_q \ge 0\}
}
\qquad\rightsquigarrow\qquad
\begin{array}{l}
\text{\upshape the premise subtrees stand unchanged;}\\
\text{\texttt{hypMultiplier\_sound} discharges the scoped row.}
\end{array}
\]
The rewrite is degenerate in the same honest sense as the base
catalogue\textquotesingle s ground-fact rule: the conclusion is not
derivable from catalogue nodes, because the base grammar has no row
form scoped to a cell-with-polynomial-atom intersection --- that scope
exists one layer up, where the child certificate that concluded
$\{g_q \ge 0\}$ lives. The two premise subtrees are ordinary
cone-family derivations, checked by the base kernel unchanged; the
node itself discharges against the verified pointwise warrant below,
which consumes exactly the two premises and nothing else.
\end{brecipe}

\begin{bprop}[semantic warrant]\label{prop:hm-sound}
If $\lambda \ge 0$ everywhere on the cell and $N + \lambda g_q \le 0$
everywhere on the cell, then $N \le 0$ everywhere on
$\mathrm{cell} \cap \{g_q \ge 0\}$.
\end{bprop}
\begin{proof}
Fix $x$ in the intersection. Then $g_q(x) \ge 0$ (membership in the
child atom) and $\lambda(x) \ge 0$ (the sign row), so
$\lambda(x)\, g_q(x) \ge 0$, and
\[
N(x) \;\le\; N(x) + \lambda(x)\, g_q(x) \;\le\; 0,
\]
the first inequality by nonnegativity of the added term, the second by
the combined row at $x$. That is the whole proof ---
\texttt{hypMultiplier\_sound} is the two displayed inequalities.
\end{proof}

\begin{bremark}[the sign row is half the rule]
Drop $\vdash \lambda \ge 0$ and the rule is unsound. Counterexample:
$N \equiv 1$, $\lambda \equiv -1$, $g_q \equiv 1$ on the cell. The
combined row holds identically ($1 - 1 = 0 \le 0$), yet $N \equiv 1$
fails everywhere on $\mathrm{cell} \cap \{g_q \ge 0\} = \mathrm{cell}$.
A free-signed multiplier converts the tax the child atom levies into a
subsidy: the combination can cancel exactly the positivity the
conclusion denies.
\end{bremark}

\begin{bremark}[the scope restriction is the other half]
The conclusion cannot claim the whole cell. With $N \equiv 1$,
$\lambda \equiv 1$, $g_q \equiv -1$, both premises hold
($\lambda \ge 0$; $1 - 1 = 0 \le 0$) while $N > 0$ everywhere --- and
the rule is nonetheless sound, because
$\mathrm{cell} \cap \{g_q \ge 0\} = \varnothing$ and the scoped claim
is vacuous. Where $g_q(x) < 0$ the term $\lambda(x) g_q(x)$ subsidizes
$N$ instead of taxing it, and the combined row bounds nothing: the
scoped conclusion is exactly as strong as the premises pay for, no
stronger.
\end{bremark}

\begin{bremark}[degree-one positivstellensatz, and where the search lives]
The move is the degree-one corner of certifying positivity over a
constrained set by nonnegative combinations: Farkas\textquotesingle{}
multiplier discipline \cite{farkas1902} (see \cite{schrijver1986} for
the sign conventions) in the affine case, sitting at the bottom of the
polynomial cone hierarchy of Krivine \cite{krivine1964} and Handelman
\cite{handelman1988} that the calculus\textquotesingle s cone rule
draws on. Nothing beyond degree one is consumed here: one multiplier,
one product, no certificate machinery beyond the two posted rows. The
\emph{search} for $\lambda$ (typically a linear program over the
cone-membership witnesses) is producer-side; only its result ships,
and the checker replays the two rows without ever seeing the search.
\end{bremark}

\begin{bremark}[the bridge hypothesis]
In the verified statement the child atom is an abstract set with one
bridge hypothesis naming its compiled sublevel reading
($\{x \mid 0 \le g_q(x)\}$). When the presentation layer of the
extension theory ships, the hypothesis discharges definitionally from
the compile; no statement changes shape. Until then it is the honest
record of exactly what the warrant assumes about the child
certificate\textquotesingle s conclusion --- its region is the
sublevel set it claims to be, and nothing more.
\end{bremark}

\begin{anchors}
Recipe \ref{rec:hm-instance} warrant & \texttt{hypMultiplier\_sound} (QtslAux/Structural.lean)\\
premise rows & cone-family sequent machinery of the base calculus (the \rulename{BernCone} and \rulename{Farkas} sections)\\
child-atom reading & the \texttt{hbridge} hypothesis of \texttt{hypMultiplier\_sound}\\
trust boundary closed & the producer-side discharge step this lemma now warrants\\
\end{anchors}
